% Источники: exam/materials/моделирование.pdf, раздел 28; https://github.com/HumsterProgrammer/iu9-notes/blob/main/8sem/Моделирование/Моделирование_лекция-04_19-02-26.md.
\section{Агрегатный подход к имитации.}

Если в системе имеет место тесное взаимодействие между компонентами, то компоненты обмениваются между собой сигналами. Каждый выходной сигнал одной компоненты является входным сигналом для другой компоненты.

Если функциональные действия аппроксимируются явно задаваемыми математическими зависимостями, позволяющими определить момент появления выходных сигналов при наличии входных сигналов, поступающих от других компонентов, то создаются условия для построения имитационной модели по модульному принципу. Каждый из таких модулей строится по унифицированной структуре и называется \textbf{агрегатом}.

\textbf{Агрегат} является математической схемой, с помощью которой возможно описание большого круга реальных процессов в системе. В любой момент $t_{ij}$ агрегат может находиться в одном из возможных состояний. Состояния агрегата $Z$ являются функциями времени.

Переход агрегата $Z$ из состояния в состояние описывается оператором перехода $H$, который позволяет по предыдущему состоянию определить очередное состояние агрегата.

Агрегат имеет входы, куда поступают входные сигналы $X(t)$ от других агрегатов, и выходы, на которых формируются выходные сигналы $Y(t)$. Кроме того, у агрегата имеются дополнительные входы, на которые поступают управляющие сигналы $g(t)$.

Выходные сигналы $Y(t)$ формируются из входных $X(t)$ и управляющих $g(t)$ сигналов оператором выхода $G$ в результате его взаимодействия с оператором $H$. Значения операторов $G$ и $H$ задаются исследователем при аппроксимации агрегатами функциональных действий реальной системы.

\clearpage
